\section{SCENARIO \#4: ''Jackson's Retreat': May 28 - June 5, 1862
}


\begin{scenariodata}
\textbf{Game Length:} 9 Turns\\
\end{scenariodata}

\begin{scenariodata}
(Weather: Turns 1-4: May 16-31 Column; Turns 5-9: June Column)\\
\textbf{First Phase:} Confederates\\
\end{scenariodata}

\begin{scenariodata}
\textbf{Union:}\\
\end{scenariodata}

\begin{scenariodata}
\textbf{Resource Factor :} 10\\
\end{scenariodata}

\begin{scenariodata}
\textbf{Initial Placement:}\\
\end{scenariodata}

\begin{scenariodata}
Hex B-37 (Franklin) : (31, 1C], (1F, 1S)\\
Hex B-25 (Moorefield) : [3t, 1A], (IF, 1S)\\
Hex A-16 (Romney) : [51, 2A, 1C], (1F, 1S)\\
Hex G-24 (Wardensville) : 251, 5A, 2C,\\
\end{scenariodata}

\begin{scenariodata}
\textbf{and adjacent :} (Fremont, 18, 1W)\\
Hex N-7 : 121, 3A, 3C, (Banks)\\
B\&O (from Hex U-14 '\\
\end{scenariodata}

to Hex Z-13): [141, 2A, 2C], (1S)

\begin{scenariodata}
B\&O (from Hex A-11 :\\
\end{scenariodata}

to Hex K-8): (31, 1A, 1C]

\begin{scenariodata}
Hex CC-33 (Warrenton) --: Available Turn \#2\\
Adds + 3 Resource Factors: 171, 5A, 5C, 1H, (Shields)\\
Hex CC-33 (Warrenton) =: Available Turn \#5\\
\end{scenariodata}

\begin{scenariodata}
Adds + 3 Resource Factors: 251,6A\\
\end{scenariodata}

\begin{scenariodata}
\textbf{Withdrawal Hexes:}\\
\end{scenariodata}

\begin{scenariodata}
N-7, 8-10, Z-13, CC-21 through CC-33\\
(Supply/Wagon Entry Hexes:\\
\end{scenariodata}

\begin{scenariodata}
N-7, S-10, Z-13, CC-21 through CC-33,\\
A-11, A-13, A-16)\\
\end{scenariodata}

\begin{scenariodata}
No additional Fortification counters,\\
\end{scenariodata}

\begin{scenariodata}
\textbf{Confederates:}\\
\end{scenariodata}

\begin{scenariodata}
Resource Factor 310\\
\end{scenariodata}

\begin{scenariodata}
\textbf{Initial Placement:}\\
\end{scenariodata}

\begin{scenariodata}
Between Hex rows M through S\$; : 421, 10A, 2C\\
South of the Potomac River; : (Jackson,\\
within 2 hexes of the Potomac. : Ewell 45, 5W)\}\\
\end{scenariodata}

\begin{scenariodata}
Hex -28 (Front Royal) : 2L, (2S, 1W\}\\
Within 8 hexes of Hex N-20\\
\end{scenariodata}

(Winchester) > 7C, 1H, (Ashby, 18, 1W)

\begin{scenariodata}
Hex N-20 (Winchester) : 21, (28, 1W)\\
Hex 0-51 (Staunton) 2 1, (1F)\\
Hex BB-53 (Charlottesville) > 11, 1A, (1F)\\
\textbf{Partisans:} McNeill, Gilmor)\}\\
\end{scenariodata}

\begin{scenariodata}
\textbf{Withdrawal Hexes:}\\
\end{scenariodata}

\begin{scenariodata}
CC-47 through CC-58\\
(Supply/Wagon Entry Hexes:\\
CC-47 through CC-58, 0-51, M-50)\\
\end{scenariodata}

\begin{scenariodata}
No additional Fortification counters.\\
\end{scenariodata}

\begin{scenariodata}
(SPECIAL \#1: Place Destruction counters on all B\&O Rail line hexes\\
\end{scenariodata}

from Hex rows L through S, including bridges. Also place Destruction counters on the bridges at N-7/N-8 and T-14/U-14.)

\begin{scenariodata}
SPECIAL \#2: In this scenario, Harper's Ferry is worth no Victory Points\\
\end{scenariodata}

to the Confederates.

\begin{scenariodata}
(SPECIAL \#3: The Confederates get 2 Victory Points for each Supply\\
\end{scenariodata}

counter they get within 15 hexes of Staunton by game's end. The Union gets 2 Victory Points for each Confederate Wagon counter on the Mapboard that is not within 15 hexes of Staunton by game's end.\}

\begin{scenariodata}
\textbf{Victory Conditions:} All other results are a draw.\\
\end{scenariodata}

\begin{scenariodata}
\textbf{ADVANCED GAME:} Union wins if they have 10 or more lead in total\\
\end{scenariodata}

\subsection*{Victory Points.}

\begin{scenariodata}
Confederates win if they have 15 or more lead in\\
\end{scenariodata}

total Victory Points.

\begin{scenariodata}
\textbf{BASIC GAME:} This scenario is not suitable for Basic Game play.\\
\end{scenariodata}

\begin{scenariodata}
\textbf{HISTORICAL RESULTS:} Confederate victory.\\
\end{scenariodata}

\begin{scenariodata}
\textbf{PLAYTESTER'S COMMENTS:} Victory in this scenario could be decided\\
\end{scenariodata}

by the weather conditions. The Confederates should mass their wagons as close to the Valley Turnpike as possible (and on the west bank of the Opequon) under heavy escort. A powerful force under Jackson (for best force-march potential) should press ahead to clear away any Union blocking forces, while Ewell brings up the rear. The key to Union success is Fremont's army. Fremont should, at all costs, try to block the Valley Turnpike to allow the other pursuers time to arrive. Fresh wagons and supplies should be brought in by the Union on Turn \#2 for both Shields and Banks.

(a)

Through a heavy rain, the Confederates moved up the Valley on the hard surface of the Valley Pike, successfully eluding the grasp of the

\begin{scenariodata}
Union pincers, floundering in the mud on all sides of them. The Union\\
\end{scenariodata}

command had to destroy Jackson's little army to justify their huge concentration of forces in the Shanandoah, but the Confederates escaped with hardly a scratch.

Successful as the Confederate maneuvers had been so far, it would not do to allow the pursuing Federals to totally overrun the Valley, so Jackson determined to make astand before withdrawing to join the main army at Richmond. With the rain-swollen Shenandoah River at flood stage, the Confederate cavalry had systematically destroyed the bridges linking the two wings of the pursuing Union forces (Fremont and Shields), and Jackson took up an excellent central position at Port Republic...

